\textbf{[To be written in step 6: policy implications of the joint within-and-across-country leakage finding.]}

The joint pattern --- both within-country across-firm leakage and across-country cross-border leakage small in Belgium under the post-2017 EU ETS price increment --- has direct implications for the design of unilateral carbon-pricing protections.

\textit{Subsections to be written:}

\subsection{Free allocation: within-country rationale empirically weaker}

\textit{The within-country rationale for free allowance allocation rests on a high elasticity of substitution between regulated and unregulated suppliers. Our $\sigma$ bound of $[0.74, 1.26]$ falls well below the elasticities the leakage-protection literature presumes. Free-allocation auctioning revenue forgone (estimated at $\sim$EUR 30bn/yr at Phase~IV prices) is hard to defend on the within-country margin alone.}

\subsection{CBAM: cross-border rationale also weaker on Belgian evidence}

\textit{The cross-border rationale for CBAM rests on observable substitution toward non-EU sources of regulated goods. The CMdG replication does not find this substitution in Belgian customs imports. CBAM's design --- which targets a small set of carbon-intensive goods --- may still be defensible on other grounds (production relocation, services, post-2019 imports) that we do not measure.}

\subsection{What we do not measure}

\textit{Three margins that our empirical strategy does not bound: (i) firm relocation outside the EU (we observe domestic Belgian firms only), (ii) substitution in services (the customs panel covers goods only), (iii) post-2019 imports (the customs panel ends at the 2020 administrative break).}

\subsection{Implications for production-network calibrations of carbon policy}

\textit{Quantitative production-network models of carbon policy \citep[e.g.,][]{baqaee2019,king2019,aghion2024} require an elasticity of substitution between regulated and unregulated suppliers as a calibration target. Our estimate provides a direct, identification-disciplined target value for one well-measured economy. The joint evidence (within $\approx$ across $\approx$ small) implies that calibrating these models on cross-border-leakage estimates from larger countries (CMdG France) would likely overstate aggregate-emissions effects of unilateral carbon pricing in dense-network small open economies.}
